\begin{notechapter}
\chapter{Activation Functions and Initialization}\label{ch:activation-functions}
\chaptersubtitle{Which nonlinearity, and how to keep signals and gradients at the right scale}

\begin{chapterintro}
The activation $\phi$ of an MLP must be nonlinear for the network to be more than a linear
model (\cref{ch:feedforward-networks}). Beyond that, its shape decides how signals and
gradients travel through many layers: whether they shrink, explode or stay usable. The
same question decides how to choose the initial weights. This chapter catalogues the usual
activations and makes precise the arguments usually given for and against them. You will
learn
\begin{itemize}[nosep]
  \item the usual activations (sigmoid, tanh, ReLU and its variants, softplus, GELU, SiLU),
        their derivatives and their ranges;
  \item what saturation and vanishing or exploding gradients are, and how the derivative of
        $\phi$ enters them;
  \item what the ``zero-centred'' argument really says, and its limits;
  \item why ReLU units can die;
  \item how to initialize the weights so that variances are preserved (LeCun, Glorot, He).
\end{itemize}
\prerequisites{\cref{ch:feedforward-networks}; the backward recursion of
\cref{backpropagation:sec:algorithm} is used in \cref{activation-functions:sec:vanishing},
and can be taken on trust at first reading.}
\end{chapterintro}

\section{What an activation should do}\label{activation-functions:sec:properties}

\begin{itemize}
  \item \textbf{Nonlinearity.} Required, and by \cref{feedforward-networks:thm:uat} any
        continuous non-polynomial $\phi$ gives universal approximation. Expressiveness is therefore
        not what distinguishes the usual choices.
  \item \textbf{Differentiability almost everywhere.} Gradient methods need $\phi'$. A finite
        number of kinks is harmless: at a kink, software uses a fixed one-sided value (for
        ReLU, $\phi'(0) = 0$), and pre-activations land exactly on a kink with probability $0$
        when they have a density.
  \item \textbf{Derivative not too small, not too large.} Backpropagation multiplies by $\phi'(z)$
        at every layer (\cref{activation-functions:sec:vanishing}). This is the property
        that matters most in deep networks.
  \item \textbf{Cheap to compute}, since it is evaluated once per unit and per example, in
        both passes.
\end{itemize}
Monotonicity, boundedness and zero-centring are often listed as well; they are neither
necessary nor sufficient (GELU and SiLU are not monotone, ReLU is neither bounded nor
centred), and their role is discussed below.

\section{A catalogue}\label{activation-functions:sec:catalogue}

\paragraph{Sigmoid and tanh.}
\begin{align*}
  \sigma(z) &= \frac1{1 + e^{-z}}, & \sigma'(z) &= \sigma(z)\bigl(1 - \sigma(z)\bigr)\leq\tfrac14,\\
  \tanh(z) &= \frac{e^z - e^{-z}}{e^z + e^{-z}} = 2\sigma(2z) - 1, & \tanh'(z) &= 1 - \tanh^2(z)\leq1 .
\end{align*}
Both are smooth, bounded and saturating: $\phi'(z)\to0$ as $\abs z\to\infty$. Because $\tanh$ is an
affine function of $\sigma(2\,\cdot)$, a tanh network and a sigmoid network represent the same functions
(\cref{activation-functions:exr:2}); they differ in how training behaves. The sigmoid
remains the output unit for binary classification (\cref{ch:sigmoid-softmax}); as a hidden
activation it has been replaced by the functions below.

\paragraph{ReLU and its variants.} The rectified linear unit $\mathrm{ReLU}(z) = \max(0, z)$ has derivative
$\bbone[z>0]$: it does not saturate for $z>0$, and its exact zeros make activations sparse. It was
popularized by Nair and Hinton (\emph{Rectified Linear Units Improve Restricted Boltzmann
Machines}, ICML, 2010) and Glorot, Bordes and Bengio (\emph{Deep Sparse Rectifier Neural
Networks}, AISTATS, 2011). For $z<0$ the derivative is $0$, which can kill units
(\cref{activation-functions:sec:dead}). Variants keep a nonzero slope there:
\begin{itemize}[nosep]
  \item \emph{leaky ReLU} $\max(\alpha z, z)$ with a small fixed $\alpha$, e.g.\ $0.01$ (Maas, Hannun and Ng,
        \emph{Rectifier Nonlinearities Improve Neural Network Acoustic Models}, ICML Workshop, 2013);
        with a learned $\alpha$ it is the \emph{PReLU} of He, Zhang, Ren and Sun (2015, below);
  \item \emph{ELU} $z$ for $z>0$ and $\alpha(e^z - 1)$ for $z\leq0$, smooth for $\alpha = 1$ and with negative outputs
        bounded below by $-\alpha$ (Clevert, Unterthiner and Hochreiter, ICLR, 2016);
  \item \emph{softplus} $\log(1 + e^z)$, a smooth ReLU: its derivative is $\sigma(z)$ and
        $0<\mathrm{softplus}(z) - \mathrm{ReLU}(z)\leq\log2$.
\end{itemize}

\paragraph{Smooth gated units.} Write the unit as the input times a gate in $[0,1]$:
ReLU is $z\cdot\bbone[z>0]$. Replacing the hard gate by a smooth one gives
\[
  \mathrm{GELU}(z) = z\,\Phi(z),\qquad\mathrm{SiLU}(z) = z\,\sigma(z),
\]
with $\Phi$ the standard normal distribution function (Hendrycks and Gimpel, \emph{Gaussian
Error Linear Units}, 2016) and $\sigma$ the sigmoid (Elfwing, Uchibe and Doya, Neural Networks,
2018; also called Swish by Ramachandran, Zoph and Le, 2017, who add a slope $z\sigma(\beta z)$). Both
are smooth, slightly negative for $z<0$ (with a minimum about $-0.17$ for GELU and $-0.28$ for
SiLU), not monotone, and close to ReLU for large $\abs z$. GELU is the default in transformers.
Since $\Phi$ has no closed form, $\mathrm{GELU}(z)\approx\frac z2\bigl(1 + \tanh\bigl(\sqrt{2/\pi}(z + 0.044715z^3)\bigr)\bigr)$ is
often used.

\begin{figure}[ht]
\centering
\begin{tikzpicture}
\begin{axis}[name=f, width=7.2cm, height=5.6cm, xmin=-4, xmax=4, ymin=-1.3, ymax=3.2,
  xlabel={$z$}, title={$\phi(z)$}, grid=major, grid style={gray!20},
  legend style={font=\small, draw=none, fill=none, at={(0.03,0.97)}, anchor=north west}, legend cell align=left]
  \addplot[very thick, bookaccent, domain=-4:4, samples=81] {1/(1+exp(-x))}; \addlegendentry{sigmoid}
  \addplot[very thick, bookexample, dashed, domain=-4:4, samples=81] {tanh(x)}; \addlegendentry{tanh}
  \addplot[very thick, bookpitfall, domain=-4:4, samples=3] {max(0,x)}; \addlegendentry{ReLU}
  \addplot[very thick, bookdef, densely dotted, domain=-4:4, samples=81] {x/(1+exp(-x))}; \addlegendentry{SiLU}
\end{axis}
\begin{axis}[at={(f.east)}, anchor=west, xshift=1.1cm, width=7.2cm, height=5.6cm, xmin=-4, xmax=4,
  ymin=-0.2, ymax=1.2, xlabel={$z$}, title={$\phi'(z)$}, grid=major, grid style={gray!20}]
  \addplot[very thick, bookaccent, domain=-4:4, samples=81] {exp(-x)/(1+exp(-x))^2};
  \addplot[very thick, bookexample, dashed, domain=-4:4, samples=81] {1-tanh(x)^2};
  \addplot[very thick, bookpitfall, jump mark left] coordinates {(-4,0) (0,0) (4,1)};
  \addplot[very thick, bookdef, densely dotted, domain=-4:4, samples=81]
    {1/(1+exp(-x))*(1+x*(1-1/(1+exp(-x))))};
\end{axis}
\end{tikzpicture}
\caption{Four activations and their derivatives. The sigmoid's derivative never exceeds $\frac14$;
tanh's reaches $1$ at $0$ only; both vanish for large $\abs z$. ReLU's derivative is $0$ or $1$. SiLU's
derivative exceeds $1$ for $z>1.28$ (by at most $0.10$, at $z\approx2.4$) and is negative for $z<-1.28$,
where SiLU decreases to its minimum.}
\label{activation-functions:fig:catalogue}
\end{figure}

\begin{table}[ht]
\centering\small
\begin{tabular}{lllll}
\toprule
& $\phi(z)$ & $\phi'(z)$ & range & $\sup\abs{\phi'}$ \\
\midrule
sigmoid & $1/(1 + e^{-z})$ & $\sigma(1 - \sigma)$ & $(0,1)$ & $\frac14$ \\
tanh & $\tanh z$ & $1 - \tanh^2$ & $(-1,1)$ & $1$ \\
ReLU & $\max(0,z)$ & $\bbone[z>0]$ & $[0,\infty)$ & $1$ \\
leaky ReLU & $\max(\alpha z, z)$ & $1$ or $\alpha$ & $\R$ & $1$ \\
ELU & $z$ or $\alpha(e^z - 1)$ & $1$ or $\alpha e^z$ & $(-\alpha,\infty)$ & $\max(1,\alpha)$ \\
softplus & $\log(1 + e^z)$ & $\sigma(z)$ & $(0,\infty)$ & $1$ \\
GELU & $z\Phi(z)$ & $\Phi(z) + z\varphi(z)$ & $[-0.17,\infty)$ & $\approx1.13$ \\
SiLU & $z\sigma(z)$ & $\sigma(z)\bigl(1 + z(1 - \sigma(z))\bigr)$ & $[-0.28,\infty)$ & $\approx1.10$ \\
\bottomrule
\end{tabular}
\caption{Usual activations; $\varphi$ is the standard normal density.}
\label{activation-functions:tab:catalogue}
\end{table}

\section{Saturation, vanishing and exploding gradients}\label{activation-functions:sec:vanishing}

\paragraph{The product structure.} Backpropagation
(\cref{backpropagation:thm:backprop}) computes the error signals $\delta^{(l)} = \partial\ell/\partial z^{(l)}$ by
\[
  \delta^{(l)} = D^{(l)}W^{(l+1)\top}\delta^{(l+1)},\qquad D^{(l)} = \diag\bigl(\phi'(z^{(l)})\bigr),
\]
so that $\delta^{(l)} = D^{(l)}W^{(l+1)\top}D^{(l+1)}W^{(l+2)\top}\cdots D^{(L-1)}W^{(L)\top}\delta^{(L)}$, and the gradient for $W^{(l)}$ is
$\delta^{(l)}a^{(l-1)\top}$. Taking operator norms,
\[
  \norm{\delta^{(l)}}\leq\Bigl(\prod_{k=l}^{L-1}\sup\abs{\phi'}\,\norm{W^{(k+1)}}\Bigr)\norm{\delta^{(L)}} .
\]
The signal reaching layer $l$ is a product of $L - l$ factors. If they are typically below $1$,
it decreases geometrically with depth: the \emph{vanishing gradient} problem, first
analysed by Hochreiter (diploma thesis, 1991) and by Bengio, Simard and Frasconi
(\emph{Learning Long-Term Dependencies with Gradient Descent is Difficult}, IEEE Transactions
on Neural Networks, 1994) for recurrent networks. If they are typically above $1$, it grows
geometrically: the \emph{exploding gradient} problem. The first layers then learn much more
slowly, or much more erratically, than the last ones.

\paragraph{The sigmoid.} With $\phi = \sigma$, every factor $\phi'$ is at most $\frac14$, so
\[
  \norm{\delta^{(l)}}\leq4^{-(L-l)}\prod_{k=l}^{L-1}\norm{W^{(k+1)}}\,\norm{\delta^{(L)}} :
\]
ten layers contribute a factor $4^{-10}\approx10^{-6}$, which the weights would have to compensate. This bound holds even when no unit is saturated.

\paragraph{Saturation.} A unit is \emph{saturated} when its pre-activation lies where $\phi$ is flat:
$\abs z$ large for sigmoid and tanh, $z<0$ for ReLU. Its derivative is then nearly $0$, and it
transmits almost no gradient, either to its incoming weights or to earlier layers. Saturation
is a property of individual units and inputs; vanishing gradients are a property of the
product along the depth. Saturation makes vanishing worse, and large weights or a poor
initialization cause saturation.

\paragraph{Remedies.} Non-saturating activations (ReLU family), an initialization that
keeps the factors near $1$ on average (\cref{activation-functions:sec:init}), normalization
layers (\cref{ch:normalization}), and residual connections $a^{(l)} = a^{(l-1)} + F(a^{(l-1)})$, whose Jacobian
$I + J_F$ does not shrink the signal, address vanishing gradients. \emph{Gradient clipping},
which rescales the gradient when its norm exceeds a threshold (Pascanu, Mikolov and Bengio,
\emph{On the Difficulty of Training Recurrent Neural Networks}, ICML, 2013), addresses
exploding gradients only: it bounds the size of a step and does nothing for gradients that
are too small.

\section{Zero-centred activations}\label{activation-functions:sec:centred}

The gradient for the incoming weights $w_j$ of unit $j$ of layer $l$ is, for one example,
\[
  \nabla_{w_j}\ell = \delta_j\,a^{(l-1)},
\]
the activation vector of the previous layer multiplied by a scalar.

\begin{proposition}\label{activation-functions:prop:orthant}
If the activations $a^{(l-1)}$ are all positive (sigmoid outputs, or the nonzero ReLU outputs),
the coordinates of $\nabla_{w_j}\ell$ for a single example all have the sign of $\delta_j$. A gradient step on
this example then increases all the weights of $w_j$, or decreases all of them.
\end{proposition}

The update of $w_j$ is confined to the positive or the negative orthant. If the weights must
move in a direction with mixed signs, say $(+1,-1)$, gradient descent on single examples can
only approximate it by alternating steps in the two allowed orthants, a zig-zag
(\cref{activation-functions:fig:orthant}). This is the argument of LeCun, Bottou, Orr and
Müller (\emph{Efficient BackProp}, 1998) for zero-mean inputs and activations, such as tanh.

\begin{figure}[ht]
\centering
\begin{tikzpicture}[scale=1.1, >=Stealth]
  \fill[bookaccent!8] (0,0) rectangle (2.4,2.0);
  \fill[bookaccent!8] (0,0) rectangle (-1.2,-1.6);
  \draw[->] (-1.4,0) -- (2.7,0) node[right] {$w_{j1}$};
  \draw[->] (0,-1.8) -- (0,2.2) node[above] {$w_{j2}$};
  \node[font=\scriptsize, bookaccent] at (1.7,1.7) {allowed};
  \node[font=\scriptsize, bookaccent] at (-0.75,-1.3) {allowed};
  \coordinate (s) at (0,0);
  \draw[->, very thick, bookdef] (s) -- (2.2,-1.1) node[right, font=\small] {target direction};
  \draw[->, thick, bookexample] (0,0) -- (1.0,0.25);
  \draw[->, thick, bookexample] (1.0,0.25) -- (0.85,-0.45);
  \draw[->, thick, bookexample] (0.85,-0.45) -- (1.85,-0.2);
  \draw[->, thick, bookexample] (1.85,-0.2) -- (1.7,-0.9);
  \node[font=\scriptsize, bookexample] at (1.15,0.55) {steps};
\end{tikzpicture}
\caption{With positive inputs, each single-example step for the weights of one unit lies in
the first or the third quadrant (shaded). Reaching a direction like $(1,-\frac12)$ takes a
zig-zag of steps of both kinds.}
\label{activation-functions:fig:orthant}
\end{figure}

\begin{pitfall}[What the argument does not say]
The constraint holds per example. A mini-batch gradient $\frac1B\sum_i\delta_j^{(i)}a^{(i)}$ sums vectors with
different signs of $\delta_j^{(i)}$ and different directions $a^{(i)}$, and is not confined to an orthant.
The effect that survives averaging is a matter of conditioning: inputs with a large common
mean make the Hessian of the unit badly conditioned (LeCun et al.\ 1998), which slows
gradient descent as in the zig-zag of \cref{gradient-tangent-space:sec:zigzag}. Centring
the inputs of every layer, which batch normalization does (\cref{ch:normalization}),
removes the problem whatever the activation. ReLU networks train well despite non-centred
activations.
\end{pitfall}

\section{Dead ReLU units}\label{activation-functions:sec:dead}

A ReLU unit is \emph{dead} if its pre-activation $w_j^\top a + b_j$ is negative for every training input. Its
output and its derivative are then $0$ on all the data, so the gradients for $w_j$ and $b_j$ vanish
on every mini-batch, and gradient descent never changes them: the unit stays dead. (Weight decay
only shrinks $w_j$ and $b_j$ towards $0$, which does not help if $b_j<0$ dominates.) A unit
switched off for a single example is not dead; that is the normal sparsity of ReLU, and
the pitfall of \cref{backpropagation:sec:pitfalls} shows its effect on one gradient.

Units die after a large step pushes $b_j$ very negative, or $w_j$ in a direction where all inputs
give negative values; high learning rates make this frequent. Leaky ReLU, ELU, GELU and SiLU
have nonzero derivatives for $z<0$ and cannot die in this sense.

\section{Initialization}\label{activation-functions:sec:init}

Weights are initialized at random, since with equal weights all units of a layer would
compute the same function and receive the same gradient forever (the \emph{symmetry} is
never broken). The scale of the random weights decides whether the signal shrinks or grows
with depth.

\paragraph{Variance propagation.} Consider $z^{(l)}_j = \sum_{k=1}^{n_{l-1}}W_{jk}a^{(l-1)}_k$ (biases initialized at $0$), with
the $W_{jk}$ independent, of mean $0$ and variance $s^2$, independent of the activations. Then
\[
  \Var\bigl(z^{(l)}_j\bigr) = n_{l-1}\,s^2\,\E\bigl[(a^{(l-1)}_k)^2\bigr],
\]
because the terms are uncorrelated with mean $0$ and $\E[W_{jk}^2a_k^2] = s^2\E[a_k^2]$. It remains to relate
$\E[a^2]$ to $\Var(z)$ for the previous layer.

\begin{itemize}
  \item \textbf{ReLU.} If $z$ is symmetric about $0$, $\E[\mathrm{ReLU}(z)^2] = \frac12\E[z^2] = \frac12\Var(z)$. Hence
        $\Var(z^{(l)}) = \frac12n_{l-1}s^2\Var(z^{(l-1)})$, and the variance is preserved if $s^2 = 2/n_{l-1}$: the
        initialization of He, Zhang, Ren and Sun (\emph{Delving Deep into Rectifiers},
        ICCV, 2015).
  \item \textbf{tanh.} Near $0$, $\tanh z\approx z$, so $\E[a^2]\approx\Var(z)$ and $s^2 = 1/n_{l-1}$ preserves the variance
        (LeCun et al.\ 1998).
  \item \textbf{Backward pass.} The same computation for $\delta^{(l)} = D^{(l)}W^{(l+1)\top}\delta^{(l+1)}$ involves $n_{l+1}$ instead of
        $n_l$. Glorot and Bengio (\emph{Understanding the Difficulty of Training Deep Feedforward
        Neural Networks}, AISTATS, 2010) compromise between the two passes with
        $s^2 = 2/(n_{l-1} + n_l)$, often with a uniform distribution on $[-\sqrt{3}s,\sqrt3s]$.
\end{itemize}

\begin{example}[Fifty layers]\label{activation-functions:ex:init}
A network of $50$ layers of width $256$, fed with $1000$ standard Gaussian inputs, with
Gaussian weights of variance $s^2$. The standard deviation of the pre-activations at several
depths:

\begin{center}
\begin{tabular}{llcccc}
\toprule
activation & $s^2$ & layer 1 & layer 10 & layer 20 & layer 50 \\
\midrule
ReLU & $1/n$ & $1.00$ & $0.044$ & $0.0011$ & $\approx0$ \\
ReLU & $2/n$ (He) & $1.41$ & $1.48$ & $1.64$ & $1.59$ \\
tanh & $1/n$ (LeCun) & $1.00$ & $0.24$ & $0.16$ & $0.10$ \\
\bottomrule
\end{tabular}
\end{center}

With $s^2 = 1/n$ and ReLU, each layer halves the variance: the prediction $2^{-(l-1)/2}$ gives
$0.044$ at layer $10$, as observed, and $4\times10^{-8}$ at layer $50$. He's initialization keeps the scale (the
first layer is $\sqrt2$ because its input is Gaussian rather than rectified). With tanh, the
scale still decays slowly, because $\abs{\tanh z}<\abs z$ shrinks every layer a little.
\end{example}

\begin{chaptersummary}
\begin{itemize}[nosep]
  \item Any continuous non-polynomial activation is expressive enough; what matters is $\phi'$, which
        multiplies the backward signal at every layer.
  \item Sigmoid ($\sigma'\leq\frac14$) and tanh saturate; ReLU has derivative $0$ or $1$; leaky ReLU, ELU, softplus,
        GELU and SiLU keep a nonzero derivative for $z<0$.
  \item $\delta^{(l)}$ is a product of $L - l$ factors $D^{(k)}W^{(k+1)\top}$: it vanishes or explodes geometrically unless the
        factors are near $1$. Clipping treats explosion only.
  \item Positive activations restrict single-example updates of a unit's weights to an orthant;
        the averaged effect is a conditioning problem, cured by centring.
  \item A ReLU unit negative on all the data gets zero gradient and stays dead.
  \item Initialize with variance $2/n_{\mathrm{in}}$ for ReLU (He), $1/n_{\mathrm{in}}$ for tanh (LeCun), or $2/(n_{\mathrm{in}} + n_{\mathrm{out}})$
        (Glorot).
\end{itemize}
\end{chaptersummary}

\section*{Exercises}
\addcontentsline{toc}{section}{Exercises}

\begin{exercise}\label{activation-functions:exr:1}
Show that $\sigma' = \sigma(1 - \sigma)$, that $\sigma'\leq\frac14$ with equality only at $0$, and that $\tanh' = 1 - \tanh^2$.
\solutionlink{activation-functions:exr:1}
\end{exercise}

\begin{exercise}\label{activation-functions:exr:2}
Show that $\tanh z = 2\sigma(2z) - 1$. Deduce that every MLP with tanh hidden units equals an MLP
with sigmoid hidden units and the same widths, and give the change of parameters.
\solutionlink{activation-functions:exr:2}
\end{exercise}

\begin{exercise}\label{activation-functions:exr:3}
Show that $\mathrm{softplus}' = \sigma$ and that $0<\mathrm{softplus}(z) - \mathrm{ReLU}(z)\leq\log2$.
\solutionlink{activation-functions:exr:3}
\end{exercise}

\begin{exercise}\label{activation-functions:exr:4}
Let $z$ have a density symmetric about $0$. Show that $\E[\mathrm{ReLU}(z)^2] = \frac12\E[z^2]$, and that
$\E[\mathrm{ReLU}(z)] = \E[\abs z]/2$, which is not $0$.
\solutionlink{activation-functions:exr:4}
\end{exercise}

\begin{exercise}\label{activation-functions:exr:5}
A ReLU unit receives inputs $a\in[0,1]^3$ and has $w = (0.5, 0.2, 0.3)$, $b = -1.2$. Show that it is dead on
every possible input. Can gradient descent, with or without weight decay, bring it back to
life? Could it if the activation were a leaky ReLU?
\solutionlink{activation-functions:exr:5}
\end{exercise}
\end{notechapter}
